\documentclass[12pt,a4paper]{article}

% ---- Encoding & Font ----
\usepackage{iftex}
\ifPDFTeX
  \usepackage[utf8]{inputenc}
  \usepackage[T1]{fontenc}
\else
  \usepackage{fontspec}  % XeTeX/LuaTeX (tectonic): Unicode nativo (·, —, ¿, ª, §)
\fi
\usepackage[spanish]{babel}
\usepackage{csquotes}

% ---- Page layout ----
\usepackage[top=2.5cm, bottom=2.5cm, left=2.5cm, right=2.5cm]{geometry}
\usepackage{setspace}
\onehalfspacing

% ---- Math ----
\usepackage{amsmath, amssymb, amsthm}
\usepackage{mathtools}
\usepackage{physics}
\usepackage{esint}  % \oiint

% ---- Graphics ----
\usepackage{graphicx}
\usepackage{tikz}
\usepackage{float}
\usepackage{caption}

% ---- Tables ----
\usepackage{array}
\usepackage{booktabs}
\usepackage{tabularx}
\usepackage{colortbl}

% ---- Colours ----
\usepackage{xcolor}
\definecolor{notebg}{HTML}{FFF3CD}
\definecolor{noteborder}{HTML}{FFC107}
\definecolor{boxred}{HTML}{F8D7DA}
\definecolor{boxredborder}{HTML}{F5C6CB}

% ---- Boxes ----
\usepackage{mdframed}
\usepackage{tcolorbox}
\tcbuselibrary{most}

\newtcolorbox{keybox}[1][]{
  colback=blue!5,
  colframe=blue!70!black,
  arc=4pt,
  boxrule=1pt,
  fonttitle=\bfseries,
  title={#1}
}

\newtcolorbox{warnbox}[1][]{
  colback=boxred,
  colframe=boxredborder,
  arc=4pt,
  boxrule=1pt,
  fonttitle=\bfseries,
  title={#1}
}

\newtcolorbox{notebox}[1][]{
  colback=notebg,
  colframe=noteborder,
  arc=4pt,
  boxrule=1pt,
  fonttitle=\bfseries,
  title={#1}
}

% ---- Hyperlinks & TOC ----
\usepackage[hidelinks]{hyperref}
\usepackage{bookmark}

% ---- Enumerate ----
\usepackage{enumitem}

% ---- Miscellaneous ----
\usepackage{icomma}
\usepackage{siunitx}
\usepackage{textcomp}
\usepackage[bottom]{footmisc}

% ---- Title ----
\title{\textbf{Clase 28: Experimento de Doble Rendija, Difracción, \\
        Coherencia e Intensidad}\\[0.3em]
        \large{Franjas de interferencia, difracción, experimento de Young
              y patrón de intensidad}}
\author{Transcripto y expandido --- Curso Física III (OpenFING)}
\date{}

\begin{document}

\maketitle
\thispagestyle{empty}
\newpage

\tableofcontents
\newpage

% ===================================================================
\section{El experimento de doble rendija}
% ===================================================================

Onda plana monocromática ($\lambda$) normal sobre dos rendijas separadas $d$;
pantalla a distancia $D$; ángulo $\theta$ hacia $P$. Dos fuentes de luz generan
oscuridad en ciertas zonas (fenómeno general de las ondas).

\subsection{Hipótesis: campo lejano y diferencia de caminos}

(1) Campo lejano $D \gg d$ (rayos casi paralelos). (2) Se ignoran aspectos
vectoriales.

\begin{keybox}[Diferencia de caminos]
  \[
  \boxed{R_2 - R_1 \approx d\,\sin\theta}
  \]
\end{keybox}

\subsection{Suma de las dos ondas}

Con amplitudes iguales y onda coherente:
\[
E = A\sin(kR_1 - \omega t) + A\sin(kR_2 - \omega t).
\]
Sumando los senos ($\Delta\phi = k(R_2-R_1)$):
\[
E = 2A\cos\!\left(\frac{k(R_2-R_1)}{2}\right)
    \sin\!\left(kR_1 - \omega t + \frac{k(R_2-R_1)}{2}\right).
\]

\subsection{Condiciones de máximo y mínimo}

\begin{keybox}[Máximos y mínimos de doble rendija]
  \[
  \boxed{d\,\sin\theta = n\,\lambda} \quad (\text{máx}), \qquad
  \boxed{d\,\sin\theta = \left(n+\tfrac12\right)\lambda} \quad (\text{mín}).
  \]
\end{keybox}

General: máximo cuando la diferencia de fases es múltiplo de $2\pi$ (diferencia
de caminos múltiplo de $\lambda$).

\subsection{Aproximación de ángulos pequeños}

Máximos y mínimos intercalados, equiespaciados en diferencia de caminos pero no
en ángulo. Para $\theta$ pequeño ($\sin\theta \approx \tan\theta = y/D$):

\begin{keybox}[Posiciones cerca del eje]
  \[
  \boxed{y_{\max} = \frac{D\lambda}{d}\,n}, \qquad
  y_{\min} = \frac{D\lambda}{d}\left(n+\tfrac12\right).
  \]
\end{keybox}

Equiespaciados con separación $D\lambda/d$; los mínimos a mitad de camino.

\subsection{Número finito de máximos}

$\sin\theta = n\lambda/d$ requiere $n\lambda/d \le 1$. Para un máximo lateral:

\begin{keybox}[Condición de interferencia]
  \[
  \boxed{\frac{d}{\lambda} \ge 1}, \qquad n_{\max} \approx \left\lfloor
  \frac{d}{\lambda}\right\rfloor.
  \]
\end{keybox}

Número finito de máximos, más espaciados hacia los lados.

% ===================================================================
\section{Difracción por una rendija}
% ===================================================================

Una rendija de ancho $a$ interfiere consigo misma. Modelándola por dos agujeros
separados $a/2$: un rayo del borde y otro del centro tienen diferencia de
caminos $\tfrac{a}{2}\sin\theta$; si es $\lambda/2$, se cancelan. Emparejando de
a dos:

\begin{keybox}[Mínimos de rendija simple]
  \[
  \boxed{a\,\sin\theta = m\,\lambda}, \qquad m \neq 0.
  \]
\end{keybox}

Los dos efectos se superponen: mínimos finos (doble rendija, $d$) y envolventes
oscuras anchas (rendija simple, $a$).

\begin{notebox}[¿Por qué un láser?]
  Es potente y muy coherente (longitud de onda definida, frente en fase en una
  región amplia). Con una lamparita el experimento es muy difícil.
\end{notebox}

% ===================================================================
\section{El experimento de Young y la coherencia}
% ===================================================================

\subsection{El experimento de Young (1801)}

Young colocó una rendija fina \emph{antes} de la doble rendija: todos sus puntos
son esencialmente el mismo, coherente consigo mismo. Así convirtió una fuente
incoherente (vela) en coherente.

\begin{warnbox}[Compromiso]
  Una rendija infinitamente fina daría coherencia perfecta pero sin intensidad.
  Young balanceó fuente grande + rendija fina.
\end{warnbox}

\subsection{Coherencia cuantitativa: el término de interferencia}

Con $E = E_1 + E_2$ y desfasaje $\phi$, la intensidad ($\propto \langle
E^2\rangle$):
\[
\frac{I}{C} = \frac{I_1}{C} + \frac{I_2}{C} + \underbrace{2\langle E_1
E_2\rangle}_{\text{interferencia}}.
\]
Si $\phi$ oscila aleatoriamente (incoherente), $\langle\cos\phi\rangle =
\langle\sin\phi\rangle = 0$:

\begin{keybox}[Fuentes incoherentes]
  \[
  \boxed{I = I_1 + I_2}
  \]
\end{keybox}

Por eso dos lamparitas (autocoherencia $\sim 10^{-8}$ s) no dan zonas oscuras.

\begin{notebox}[Coherente]
  Con fase fija, el término de interferencia sobrevive y aparecen las franjas.
  Importan la coherencia temporal y la espacial.
\end{notebox}

% ===================================================================
\section{Cálculo de la intensidad}
% ===================================================================

Con rendijas muy finas (difracción ignorada), cada campo decae como $1/R$:
\[
E = \frac{A}{R_1}\sin(kR_1 - \omega t) + \frac{A}{R_2}\sin(kR_2 - \omega t).
\]
Campo lejano: $R_2 \approx R_1$ en el prefactor, pero no en la fase. Sumando y
promediando en el tiempo ($\langle\sin^2\rangle = 1/2$):
\[
I = \frac{2A^2}{R_1^2}\cos^2\!\left(\frac{k(R_2-R_1)}{2}\right).
\]
Normalizando por una rendija ($I_1 = CA^2/2R_1^2$):

\begin{keybox}[Intensidad de doble rendija]
  \[
  \boxed{\frac{I}{I_1} = 4\cos^2\!\left(\frac{k(R_2 - R_1)}{2}\right)}
  = 4\cos^2\!\left(\frac{k\,d\,\sin\theta}{2}\right).
  \]
\end{keybox}

\begin{itemize}[nosep]
  \item El coseno se anula en los mínimos (zonas oscuras).
  \item En los máximos la intensidad es el \textbf{cuádruple} (no el doble): la
        amplitud se duplica y la intensidad va con su cuadrado.
  \item Hay una dependencia suave extra $1/R_1^2$ (decae hacia los bordes),
        distinta de la interferencia.
\end{itemize}

\begin{notebox}[Fin del curso]
  Con el cálculo de la intensidad del experimento de doble rendija cierra el
  curso de Física III.
\end{notebox}

\end{document}
